\section{Praktikum}
\label{sec:praktikum-bab2}

\begin{praktikumbox}
\textbf{Judul:} Konversi bilangan dan verifikasi penyandian data di
memori.\\[0.3em]
\textbf{Sub-CPMK:} fondasi untuk 1.1--4.3 \quad\textbullet\quad
\textbf{Alokasi waktu:} 100 menit \quad\textbullet\quad
\textbf{Bentuk:} individu
\end{praktikumbox}

\subsection{Tujuan Praktikum}

Memastikan mahasiswa dapat melakukan konversi bilangan secara manual
dan memverifikasi sendiri bahwa penyandian data di memori sesuai
dengan yang dihitung.

\subsection{Alat dan Bahan}

\begin{itemize}
    \item Komputer dengan Emu8086 dan DOSBox terpasang.
    \item Program pada Kode~\ref{lst:penyandian} dan
          Kode~\ref{lst:flag-uji}.
    \item Kalkulator \eng{programmer} (boleh kalkulator bawaan sistem
          operasi) untuk memeriksa hasil, bukan untuk menggantikan
          perhitungan manual.
\end{itemize}

\subsection{Langkah Kerja}

\begin{enumerate}
    \item \textbf{Konversi manual.}
    \begin{enumerate}
        \item Konversi desimal \dec{200}, \dec{100}, dan \dec{63} ke
              biner dan heksadesimal tanpa alat bantu.
        \item Nyatakan \dec{-1}, \dec{-64}, dan \dec{-100} dalam
              \eng{two's complement} 8 bit.
        \item Periksa hasil Anda dengan kalkulator \eng{programmer}.
              Bila berbeda, temukan letak kesalahannya.
    \end{enumerate}

    \item \textbf{Verifikasi penyandian di memori.}
    \begin{enumerate}
        \item Rakit dan jalankan program pada Kode~\ref{lst:penyandian}
              di Emu8086.
        \item Buka jendela memori dan arahkan ke segmen data.
        \item Isi tabel pengamatan: untuk setiap variabel, catat
              alamat dan isi byte-nya.
        \item Bandingkan hasil pengamatan dengan perhitungan manual
              Anda pada Tabel~\ref{tab:penyandian-hasil}.
    \end{enumerate}

    \item \textbf{Menguji pengaruh flag.}
    \begin{enumerate}
        \item Jalankan program pada Kode~\ref{lst:flag-uji} dengan
              \eng{single step}.
        \item Sebelum eksekusi \asm{ADD AL, 01h}, tuliskan prediksi
              Anda untuk \flag{CF}, \flag{ZF}, \flag{SF}, dan
              \flag{OF}.
        \item Setelah eksekusi, catat nilai sebenarnya dan bandingkan
              dengan prediksi.
        \item Ulangi dengan nilai \asm{7Fh} serta \asm{80h} dan
              jelaskan perbedaan hasilnya.
    \end{enumerate}

    \item \textbf{Menyusun tabel konversi pribadi.}
    \begin{enumerate}
        \item Buat tabel heksadesimal 0--F beserta nilai biner dan
              desimalnya di buku catatan Anda.
        \item Hafalkan tabel tersebut. Tabel ini akan sangat
              mempercepat pekerjaan Anda pada bab-bab berikutnya.
    \end{enumerate}
\end{enumerate}

\subsection{Hasil yang Dikumpulkan}

\begin{enumerate}
    \item Lembar kerja konversi manual (sebelum diperiksa dengan
          kalkulator).
    \item Tabel pengamatan alamat dan isi memori.
    \item Tabel prediksi versus hasil pengamatan flag.
    \item Tabel konversi heksadesimal pribadi.
\end{enumerate}

\begin{asesmenbox}
\textbf{Asesmen fondasi (prasyarat, tanpa nilai numerik).} Praktikum
ini tidak dinilai, tetapi menjadi prasyarat untuk Tugas 1 pada
Bab~\ref{chap:register}. Mahasiswa yang belum menguasai konversi
bilangan akan mengalami kesulitan besar pada bab-bab berikutnya.

\begin{tabular}{@{}p{7.0cm}p{6.4cm}@{}}
    \toprule
    \textbf{Butir} & \textbf{Kriteria terpenuhi} \\
    \midrule
    Konversi manual & Seluruh soal konversi benar dan langkah
    pengerjaan terlihat \\
    \addlinespace
    Verifikasi memori & Isi memori hasil pengamatan sama dengan
    perhitungan manual \\
    \addlinespace
    Prediksi flag & Minimal tiga dari empat prediksi benar \\
    \addlinespace
    Tabel pribadi & Tabel 0--F lengkap dan sudah dihafalkan \\
    \bottomrule
\end{tabular}
\end{asesmenbox}
